\section{Interface index: what stands between the corpus and a proof}

This section is the join layer of the paper. Parts I--II inventory what the corpus proves; this
part inventories, obligation by obligation, exactly what is \emph{missing} to turn each proved
reduction into a proof of Erd\H{o}s \#249, in a form usable as a premise without opening the Lean
tree. Erd\H{o}s \#249 asks whether $S = \sum_{n\ge 0} \varphi(n)/2^n$ is irrational. It is OPEN.
Nothing in this section decides it. Every row below either (a) states a proved theorem exactly, or
(b) states an unproved but well-posed proposition (an \emph{open producer} or \emph{open supply})
together with the exact proved consumer that would turn it into irrationality of $S$.

\subsection{Reading this index}

Every entry follows the same discipline. \textbf{Yields}: the exact statement proved, with
quantifiers in their proved order. \textbf{Exact mismatch}: precisely where the proved statement
falls short of the open obligation it is compared against --- never ``it doesn't work,'' always the
named axis of shortfall. \textbf{What would close it}: the exact remaining proposition, stated in
full, whose proof (via the cited Lean consumer, already on disk) proves $\mathrm{Irrational}(S)$.

Each numbered claim carries three tags. \scale{fixed} means proved at finitely many explicit
numerals with no argument in the growing parameter; \scale{bounded} means proved on a bounded range
of a parameter that is itself unbounded elsewhere in the statement; \scale{uniform} means proved for
\emph{every} value of the relevant parameter (the strongest positive tag short of matching the
open target); \scale{cofinal} means the statement itself has the shape \m{\forall a_0\ \exists a\ge
a_0}, i.e.\ it is exactly the quantifier shape the open obligation needs; \scale{n/a} applies to
structural non-existence results with no growing parameter. The evidence tags are
\ev{Lean}, \ev{Cert}, \ev{Math}, \ev{Cited}, \ev{Open}, used exactly as defined in Part I.

The canonical open obligation for \#249, proved in Part II to be sufficient
(\lean{irrational_totient_series_of_certificate_supply}{TotientTailPeriodKiller.lean:394}), is the
\textbf{249-supply} obligation:
\[
  \forall h \ge 1,\ \forall N_0,\ \exists N \ge N_0,\ \exists L,\ \mathtt{certifiedKill}\ h\ N\ L .
\]
Fourteen near-miss rows are catalogued against it below: one headline row (the sharpest reduction
in the corpus, given its own subsection) and thirteen further rows, each attacking the same
obligation from a different coordinate.

\subsection{The headline: \#249 reduces to a single unsupplied exponential-sum bound}
\label{ssec:headline}

The following two definitions, verified against the live Lean tree in this session, are the
sharpest statement of the \#249 frontier that the corpus contains.

\begin{defn}[Window discrepancy]
\lean{windowDiscrepancy}{TotientTailPeriodKiller.lean:66}. For $h,N,L\in\mathbb N$,
\[
  A_{h,N,L} \;=\; \sum_{j=0}^{L-1} \bigl(\varphi(N+h+1+j) - \varphi(N+1+j)\bigr)\cdot 2^{\,L-1-j}
  \ \in \mathbb Z .
\]
This is the depth-$L$ truncation of $2^L\cdot(R_{N+h}-R_N)$, where
$R_N = \sum_{j\ge 1}\varphi(N+j)/2^j$ is the local totient tail
(\lean{totientTail}{TotientTailPeriodKiller.lean:57}).
\end{defn}

\begin{defn}[Certified kill]
\lean{certifiedKill}{TotientTailPeriodKiller.lean:72}.
\[
  \mathtt{certifiedKill}\ h\ N\ L \;:\Leftrightarrow\;
  (N{+}h{+}L{+}2) < A_{h,N,L}\bmod 2^L < 2^L - (N{+}h{+}L{+}2).
\]
That is: the residue of $A_{h,N,L}$ modulo $2^L$ avoids the radius-$(N{+}h{+}L{+}2)$ neighbourhood
of $0$. \lean{certifiedKill_depth_floor}{TotientTailPeriodKiller.lean:79} records the necessary
room condition this forces: $2(N{+}h{+}L{+}2) < 2^L$.
\end{defn}

\begin{defn}[First-harmonic real and complex characters]
\lean{windowFirstCos}{FirstHarmonicGap.lean:27} and
\lean{windowFirstExp}{FirstHarmonicPivot.lean:38}:
\[
  \mathtt{windowFirstCos}\ h\ N\ L = \cos\!\Bigl(2\pi\,\tfrac{A_{h,N,L}\bmod 2^L}{2^L}\Bigr),\qquad
  \mathtt{windowFirstExp}\ h\ N\ L = \exp\!\Bigl(i\,2\pi\,\tfrac{A_{h,N,L}\bmod 2^L}{2^L}\Bigr).
\]
\lean{windowFirstExp_re}{FirstHarmonicPivot.lean:42} records $\mathrm{Re}(\mathtt{windowFirstExp})
= \mathtt{windowFirstCos}$ and \lean{norm_windowFirstExp}{FirstHarmonicPivot.lean:48} that
$\|\mathtt{windowFirstExp}\,h\,N\,L\| = 1$: this is literally the first additive character of the
endpoint discrepancy modulo $2^L$, a standard exponential-sum object.
\end{defn}

\begin{thm}[Real-part certificate consumer]
\label{thm:hgap-real}
\lean{exists_certifiedKill_of_first_harmonic_gap}{FirstHarmonicGap.lean:128}. \ev{Lean}.
\scale{uniform}. \coord{first-harmonic}.
For all $h,X,L$ with $0<X$ and the room condition $16(2X{+}h{+}L{+}2)\le 2^L$, if
\[
  \sum_{N=X}^{2X-1} \mathtt{windowFirstCos}\ h\ N\ L \;\le\; \tfrac{9}{10}\,X ,
\]
then $\exists N\in[X,2X)$ with $\mathtt{certifiedKill}\ h\ N\ L$.
\end{thm}

\begin{thm}[Subset consumer --- no density or partition hypothesis]
\label{thm:hgap-subset}
\lean{exists_certifiedKill_of_first_harmonic_gap_subset}{FirstHarmonicGap.lean:104}.
\ev{Lean}. \scale{uniform}. \coord{first-harmonic}.
For any nonempty finite $T\subseteq\mathbb N$ with $T\subset[0,2X)$ and the same room condition, if
\[
  \sum_{N\in T} \mathtt{windowFirstCos}\ h\ N\ L \;\le\; \tfrac{9}{10}\,|T| ,
\]
then $\exists N\in T$ with $\mathtt{certifiedKill}\ h\ N\ L$. This is strictly stronger than
Theorem~\ref{thm:hgap-real}: $T$ can be \emph{any} explicitly chosen nonempty finite subset of the
dyadic block, not the whole block, and the proof (an averaging pigeonhole, \S below) never uses
that $T$ has positive density or comes from a partition.
\end{thm}

\begin{thm}[Complex norm consumer and the open producer]
\label{thm:hgap-norm}
\lean{exists_certifiedKill_of_first_harmonic_norm_gap}{FirstHarmonicPivot.lean:53} and
\lean{irrational_totient_series_of_first_harmonic_norm_gap}{FirstHarmonicPivot.lean:83}.
\ev{Lean} for both implications; the hypothesis \m{DTWFirstHarmonicNormGap} is \ev{Open}.
The complex norm bound is strictly stronger than the real-part bound
(\m{|z|\ge\mathrm{Re}(z)}, and $21/25 < 9/10$ absorbs the slack), so it composes through
Theorem~\ref{thm:hgap-real} to the same certificate. Define
\[
  \mathtt{DTWFirstHarmonicNormGap} :\Leftrightarrow\;
  \forall h{>}0\ \forall X_0\ \exists X, L,\ \max(X_0,1)\le X\ \wedge\ 16(2X{+}h{+}L{+}2)\le 2^L\ \wedge
\]
\[
  \Bigl\|\ \sum_{N=X}^{2X-1}\mathtt{windowFirstExp}\ h\ N\ L\ \Bigr\| \;\le\; \tfrac{21}{25}\,X .
\]
\coord{first-harmonic}. \scale{cofinal} (this is the open target itself). Then
\[
  \mathtt{DTWFirstHarmonicNormGap} \;\Longrightarrow\; \mathrm{Irrational}\Bigl(\sum_{n\ge 0}
  \tfrac{\varphi(n)}{2^n}\Bigr),
\]
proved in full, with no further gap, by chaining
Theorem~\ref{thm:hgap-norm}$\to$Theorem~\ref{thm:hgap-real}$\to$
\lean{irrational_totient_series_of_certificate_supply}{TotientTailPeriodKiller.lean:394}.
\end{thm}

\begin{obs}[The exact negative]
\label{obs:no-instance}
No instance of \m{\mathtt{DTWFirstHarmonicNormGap}}'s hypothesis, nor of \m{hgap} in
Theorems~\ref{thm:hgap-real}--\ref{thm:hgap-subset}, is proved anywhere in the corpus, at any $h$,
$X$, $L$. \ev{Open} is exact, not conservative rounding. The two consumer theorems that
\emph{use} the subset form (Theorem~\ref{thm:hgap-subset}) are conditional consumers, not a
supply: \lean{AdjacentPhaseSeparation.lean}{AdjacentPhaseSeparation.lean:267} instantiates $T$ as
the explicit two-element set $\{N,M\}$ (\texttt{apply exists\_certifiedKill\_of\_first\_harmonic\_gap\_subset (\{N, M\} : Finset Nat)}). The two call sites in
\lean{PivotAntiReconstruction.lean}{PivotAntiReconstruction.lean:882} and
\lean{PivotAntiReconstruction.lean}{PivotAntiReconstruction.lean:1572} use, respectively, an
arbitrary supplied finite set $T$ and the explicit singleton $\{N\}$; neither call site supplies
such a set or discharges its harmonic-gap hypothesis. There is no theorem anywhere in the corpus
that discharges $hgap$ at a single instance, let alone cofinally.
\end{obs}

\begin{rem}[Why this is the sharpest reduction, and why it is tractable]
This reduction converts the open producer for \#249 from certificate-supply language --- ``exhibit
a certified kill $(N,L)$'' --- into a \emph{constant-saving cancellation estimate for the first
additive character of the totient window discrepancy over a dyadic block}: exactly the shape of a
classical exponential-sum bound in analytic number theory (a Weyl-type estimate for
$\sum_N e(\theta_N)$). Three features make this genuinely the narrowest gap in the corpus.

First, the saving required is a \emph{constant}, not $o(X)$: $21/25$ (complex norm) or $9/10$
(real part), fixed independent of $h$, $X$, $L$. Any nontrivial cancellation estimate --- even one
far weaker than square-root cancellation --- would suffice; this is not asking for GRH-strength
input.

Second, the room condition $16(2X{+}h{+}L{+}2)\le 2^L$ forces only $L \gtrsim \log_2 X + 5$, which
is satisfiable for every $X$ by simply taking $L$ large enough (the depth floor grows logarithmically
in $X$, not polynomially); no scale obstruction of the kind documented elsewhere in this paper
(e.g.\ the $\sqrt{2\log_2 N}$ silent-position defect for \#257, Part~I) stands in the way of the
producer's own room requirement.

Third, the subset form (Theorem~\ref{thm:hgap-subset}) removes even the requirement that the
cancellation happen on the whole dyadic block or a positive-density subset of it: any single
explicitly named finite $T$ with the averaged bound would do. In particular a supplier-fibre or
sparse-subsequence argument --- e.g.\ picking $N$ to be one less than a prime, as the exact pivot
algebra in \lean{FirstHarmonicPivot.lean}{FirstHarmonicPivot.lean} already sets up via
\texttt{pivotArgument}, \texttt{pivotPrime}, \texttt{pivotSupplierBases} --- is a legitimate route
to this obligation and does not need to control the discrepancy on a full interval.
\end{rem}

\subsection{Row-by-row: the remaining near misses against 249-supply}

The headline (\S\ref{ssec:headline}) is row \textbf{e1-companion} of the source interface catalogue
in full. The remaining thirteen rows below attack the same 249-supply obligation from independent
coordinates in the Lean tree: the actual-LCM sign corridor, the top-edge staircase, the raw
rational approximant, the short-window arithmetic word, the diagonal pincer certificate bank, the
$h$-uniform single-window certificate, the Farey/continued-fraction denominator growth law, the
2-adic pulse block, the LCM-jump slack scalar, the prime-jump commutator, the M\"obius--Mersenne
denominator channel, and the carry-kernel rank. None of them is closed; each is recorded with its
exact remaining content.

\subsubsection{SGN-01 --- the positive-sign half of the certified-kill band (\coord{actual-lcm-sign})}

\begin{defn}
The \emph{actual LCM tail orbit} at exponent $a$ is
$\mathtt{actualLcmTailOrbit}\ a = R_{2H+2H} - R_{H+... }$, precisely
$\mathrm{totientTail}(2H)-\mathrm{totientTail}(H)$ where $H = \mathtt{periodLcm}(2^a)$
(\lean{actualLcmTailOrbit_pos}{TotientActualLcmOrbitSign.lean:144}).
\end{defn}

\begin{thm}
\lean{actualLcmTailDiff_shift_pos}{TotientActualLcmOrbitSign.lean:39} and
\lean{actualLcmTailOrbit_pos}{TotientActualLcmOrbitSign.lean:144}. \ev{Lean}. \scale{uniform}.
For every $a\ge 8$ and every $J$ with $J+(a{+}6) < 2\cdot 2^a$,
\[
  0 \;<\; \mathrm{totientTail}(2H{+}J) - \mathrm{totientTail}(H{+}J), \qquad H=\mathtt{periodLcm}(2^a).
\]
Hence $0 < \mathtt{actualLcmTailOrbit}\ a$ for every $a\ge 8$. The proof is genuinely uniform in
$a$: it rests on two facts proved for all $a\ge 8$ by a two-case structural split (divisor letter
vs.\ foreign prime power), with absolute constants $4,8,32$ and no lookup table ---
\lean{lcmRayArithmeticLetter_pos_of_lt_two_mul}{TotientActualLcmOrbitArithmetic.lean:1703} and
\lean{periodLcm_lt_eight_mul_t_mul_lcmRayArithmeticLetter}{TotientActualLcmOrbitArithmetic.lean:1906}.
\end{thm}

\begin{rem}[Exactly half the certificate band, and where the other half goes]
$\mathtt{certifiedKill}$ needs the residue to avoid a \emph{symmetric} radius-$(N{+}h{+}L{+}2)$
neighbourhood of $0$ on both sides. SGN-01 gives positivity only. Its companion
\lean{actualLcm_integral_forces_topEdgeResidue}{TotientActualLcmOrbitSign.lean:211} proves that,
under integrality, the consequence is not a kill but the opposite: the residue is forced to the
\emph{top edge}, exactly $2^K - e$ where $e$ is the true carry orbit (\lean{carryOrbit}{CarrySurvivorExtinction.lean:387}). So the
lower half of the certified-kill band is discharged unconditionally and cofinally by SGN-01; the
upper half is untouched, and what remains open is a genuinely different, one-sided statement about
the top edge --- row TE-04 below. The corpus's own normal-form tables tag this row
\scale{bounded}; that tag is misleading, since only the window offset $J$ is bounded relative to
the height, and the height $2\cdot 2^a$ itself is unbounded: the correct tag is \scale{uniform}.
\end{rem}

\subsubsection{TE-04 --- the one-sided top-edge staircase (\coord{actual-lcm-top-edge})}

This is the sharpest \emph{landed weakening} of the 249-supply obligation in the corpus: the
symmetric certified-kill band is replaced by a single upper inequality, halving the information
the open producer must supply.

\begin{thm}[Consumers, fully proved]
\lean{actualLcmTailDiff_notMem_int_of_topEdgeResidueGap}{TotientActualLcmTopEdgeStaircase.lean:1260},
\lean{actualLcmTailOrbit_notMem_int_of_topEdgeResidueGap}{TotientActualLcmTopEdgeStaircase.lean:1314},
\lean{irrational_totientSeries_of_topEdgeResidueGapSupply}{TotientActualLcmTopEdgeStaircase.lean:2184}.
\ev{Lean}. \scale{uniform}. For every $a\ge 8$ and every $J,K,m$ inside the sign corridor
($J{+}K{+}(a{+}6) < 2\cdot 2^a$), a one-sided residue gap at precision $m\le K$ ---
room $2H{+}J{+}K{+}2 < 2^m$ and $A_{H,H+J,K}\bmod 2^m \le 2^m - (2H{+}J{+}K{+}2)$ --- already
forces $\mathrm{totientTail}(2H{+}J) - \mathrm{totientTail}(H{+}J)\notin\mathbb Z$. No lower margin
at all is demanded. The proof chain is complete: the theorem holds for every $a\ge 8$, and
\emph{\m{\mathtt{PowerTwoActualLcmTopEdgeResidueGapSupply}} $\Rightarrow$ Irrational $S$} is proved.
\end{thm}

\begin{defn}[The open producer]
\lean{PowerTwoActualLcmTopEdgeResidueGapSupply}{TotientActualLcmTopEdgeStaircase.lean:1325}.
\ev{Open}. \scale{cofinal}.
\[
  \forall a_0\ \exists a,K,m,\ a_0\le a\ \wedge\ 8\le a\ \wedge\ K{+}(a{+}6)<2\cdot 2^a\ \wedge\
  \mathtt{ActualLcmTopEdgeResidueGap}\ a\ 0\ K\ m .
\]
No unconditional statement anywhere in the corpus locates the residue of the diagonal word
$A_{H,H,K}$ (where $H=\mathtt{periodLcm}(2^a)$) below the top strip at even one large $a$.
\lean{actualLcmTopEdgeResidueGap_iff_terminal}{TotientActualLcmTopEdgeStaircase.lean:909} further
reduces the condition to the last $m$ arithmetic letters of the word alone.
\end{defn}

\begin{prop}[The five strictly weaker sufficient links --- proving any one closes \#249]
\label{prop:te-chain}
All five are proved sufficient for
\lean{PowerTwoActualLcmTopEdgeResidueGapSupply}{TotientActualLcmTopEdgeStaircase.lean:1325} in the
same module, and each is \emph{strictly weaker} in the sense that it drops a hypothesis, widens a
band, or asks only for magnitude rather than a signed inequality.
\begin{enumerate}[leftmargin=*]
  \item \lean{PowerTwoAdjacentSuffixMidbandSupply}{TotientActualLcmTopEdgeStaircase.lean:1334}
    (\ev{Open}, \scale{cofinal}):
    \[
      \forall a_0\ \exists a,m,\ a_0\le a \wedge 8\le a \wedge m{+}1{+}(a{+}6)<2\cdot 2^a \wedge
      2H{+}m{+}3 < 2^m \wedge
    \]
    \[
      2H{+}m{+}2 \le \mathtt{diagonalAdjacentSuffixResidue}(2^a)\,0\,m \le 2^m - (2H{+}m{+}2)
    \]
    (a two-sided band on the adjacent-suffix residue directly, one candidate depth $m$).
  \item \lean{PowerTwoOddGuardTopEdgeHalfWordBandSupply}{TotientActualLcmTopEdgeStaircase.lean:1349}
    (\ev{Open}, \scale{cofinal}): at the odd guarded depth $2q{+}1$, a two-sided band of half-width
    $H{+}q{+}2$ on the half-word residue mod $4^q$ --- ``substantially weaker than the older fixed
    $1/32$ central band.''
  \item \lean{PowerTwoActualFinalTopEdgeMagnitudeSupply}{TotientActualLcmTopEdgeStaircase.lean:1471}
    (\ev{Open}, \scale{cofinal}), \emph{proved equivalent} to the previous one via
    \lean{topEdgeHalfWordBandSupply_iff_actualFinalCenteredMagnitudeSupply}{TotientActualLcmTopEdgeStaircase.lean:1479}:
    \[
      \forall a_0\ \exists a,q,\ \max(14,a_0)\le a \wedge \mathtt{oddGuardedCanonicalAdjacentSuffixDepth}(2^a)=2q{+}1
      \wedge H{+}q{+}2 \le |\mathtt{actualOddHalfCenteredLift}\ a\ q| .
    \]
  \item \lean{PowerTwoFlexibleActualTopEdgeMagnitudeSupply}{TotientActualLcmTopEdgeStaircase.lean:1927}
    (\ev{Open}, \scale{cofinal}): the same magnitude bound with the depth restriction relaxed from
    ``canonical guarded'' to any odd $2q{+}1$ satisfying the half-cell fit
    $2(H{+}q{+}2)\le 4^q$ and the sign-corridor room $2q{+}2{+}(a{+}6)<2\cdot 2^a$.
  \item \lean{PowerTwoFlexibleActualTerminalDominanceSupply}{TotientActualLcmTopEdgeStaircase.lean:1908}
    --- the weakest of all five; treated separately as row TE-05-weakest below.
\end{enumerate}
Proving \emph{any one} of these five closes Erd\H{o}s \#249; the corpus has already proved all the
implications from each to
\lean{PowerTwoActualLcmTopEdgeResidueGapSupply}{TotientActualLcmTopEdgeStaircase.lean:1325} (chain:
midband $\to$ residue-gap at \lean{}{TotientActualLcmTopEdgeStaircase.lean:2109}; half-word-band
$\to$ midband at \lean{}{TotientActualLcmTopEdgeStaircase.lean:2058}; final-magnitude
$\Leftrightarrow$ half-word-band at \lean{}{TotientActualLcmTopEdgeStaircase.lean:1479};
flexible-magnitude $\to$ midband at \lean{}{TotientActualLcmTopEdgeStaircase.lean:2008};
final-magnitude $\to$ flexible-magnitude at \lean{}{TotientActualLcmTopEdgeStaircase.lean:1993}).
\end{prop}

\subsubsection{TE-05-weakest --- one scalar inequality at cofinally many odd ranks (\coord{actual-lcm-top-edge})}

This is the weakest landed link in the whole \#249 sufficiency lattice: a single comparison
between one terminal arithmetic letter and twice a centred lift.

\begin{defn}
\lean{PowerTwoFlexibleActualTerminalDominanceSupply}{TotientActualLcmTopEdgeStaircase.lean:1908}.
\ev{Open}. \scale{cofinal}. Writing $H=\mathtt{periodLcm}(2^a)$,
\[
  \forall a_0\ \exists a,q,\ a_0\le a \wedge 8\le a \wedge 2q{+}2{+}(a{+}6)<2\cdot 2^a \wedge
  2(H{+}q{+}2)\le 4^q \wedge
\]
\[
  \mathtt{diagonalWindowIncrement}(2^a)(2q{+}2) \;\le\; 2\cdot \mathtt{actualOddHalfCenteredLift}\ a\ q .
\]
\end{defn}

\begin{thm}[Consumer]
\lean{actualLcmTailOrbit_notMem_int_of_actualTerminalDominance}{TotientActualLcmTopEdgeStaircase.lean:1880}
and
\lean{irrational_totientSeries_of_flexibleActualTerminalDominanceSupply}{TotientActualLcmTopEdgeStaircase.lean:2221}.
\ev{Lean}. \scale{uniform}. The dominance hypothesis at a single odd rank already excludes
integrality of \m{\mathtt{actualLcmTailOrbit}\ a}, and the supply predicate composes to
\m{\mathrm{Irrational}(S)}.
\end{thm}

\begin{obs}[The exact identity, and the internal tension it exposes]
\label{obs:staircase-tension}
\lean{two_mul_actualOddHalfCenteredLift_eq_terminal_sub_trueCarry}{TotientActualLcmTopEdgeStaircase.lean:1678}.
\ev{Lean}. \scale{uniform}. For $a\ge 8$, room $2q{+}1{+}1{+}(a{+}6)<2\cdot2^a$, half-cell fit
$2(H{+}q{+}2)\le 4^q$, and any integral representative $z$ with $z = \mathtt{actualLcmTailOrbit}\ a$:
\[
  2\cdot \mathtt{actualOddHalfCenteredLift}\ a\ q
  \;=\; \mathtt{diagonalWindowIncrement}(2^a)(2q{+}2) - \mathtt{carryOrbit}\ H\ H\ z\ (2q{+}1) .
\]
This is an \emph{exact equality}, not a bound. Comparing it against the dominance inequality above,
\emph{the dominance inequality is literally equivalent to
$\mathtt{carryOrbit}\,H\,H\,z\,(2q{+}1)\le 0$}. But SGN-02
(\lean{actualLcm_trueEndpointSurvivor_neg}{TotientActualLcmOrbitSign.lean:172}, \ev{Lean},
\scale{uniform}) proves that under integrality the true carry orbit is \emph{strictly positive}
throughout this exact corridor: for $a\ge 8$ and $J{+}K{+}(a{+}6)<2\cdot 2^a$,
\[
  0 < \mathtt{carryOrbit}\ H\ (H{+}J)\ d\ K
\]
whenever $d$ is the real-valued representative of the translated tail difference. So the two
branches of the corridor-escape disjunction that
\lean{actualLcmTailOrbit_notMem_int_of_actualTerminalCarryCorridorEscape}{TotientActualLcmTopEdgeStaircase.lean:1824}
offers are \emph{not} symmetric: SGN-02 has already eliminated the branch that the identity most
naturally supplies (dominance, $\mathtt{carryOrbit}\le 0$), and the surviving branch is the
\emph{lower} escape
\[
  2\cdot \mathtt{actualOddHalfCenteredLift}\ a\ q \;\le\;
  \mathtt{diagonalWindowIncrement}(2^a)(2q{+}2) - (2H{+}2q{+}3) .
\]
This is stated precisely so the reader can check it: it is neither a refutation of the dominance
route (SGN-02 constrains the sign of $\mathtt{carryOrbit}$, it does not itself bound
$\mathtt{diagonalWindowIncrement}$ or $\mathtt{actualOddHalfCenteredLift}$ against $0$, so dominance
is not shown \emph{false}, only shown to entail a carry-orbit sign the census never violates) nor a
proof (no theorem excludes dominance outright). It is the precise reason the finite census keeps
landing just inside the corridor rather than outside it, and it identifies the lower-escape branch
as the coordinate where a producer is actually needed.
\end{obs}

\begin{prop}[What would close it]
Either (i) the lower-escape branch cofinally --- as displayed just above --- or (ii) the two-sided
magnitude form \lean{PowerTwoFlexibleActualTopEdgeMagnitudeSupply}{TotientActualLcmTopEdgeStaircase.lean:1927}
(item 4 of Proposition~\ref{prop:te-chain}), which asks only $H{+}q{+}2 \le
|\mathtt{actualOddHalfCenteredLift}\ a\ q|$ and which the file proves
(\lean{flexibleActualTerminalCarryCorridorEscapeSupply_of_magnitude}{TotientActualLcmTopEdgeStaircase.lean:1938})
implies corridor escape via a clean sign split (positive branch escapes above the terminal letter,
negative branch escapes below the directed bound), sidestepping the SGN-02 tension entirely because
it does not commit to a sign for the centred lift.
\end{prop}

\subsubsection{SEP-02 --- separation from an explicit rational approximant (\coord{raw-approximant})}

\begin{thm}
\lean{abs_actualLcmTailOrbit_sub_rawApprox_lt}{TotientActualLcmOrbitSeparation.lean:141} and
\lean{irrational_totientSeries_of_actualLcmOrbitSeparationSupply}{TotientActualLcmOrbitSeparation.lean:305}.
\ev{Lean}. \scale{uniform}. Unconditionally, for every $a$ and $q$,
\[
  \bigl|\, \mathtt{actualLcmTailOrbit}\ a - \mathtt{actualLcmRawApprox}\ a\ q \,\bigr|
  \;<\; \frac{4H + 2(2q{+}1) + 4}{2^{2q+2}} ,
\]
where $\mathtt{actualLcmRawApprox}\ a\ q$ is an explicit finite computable rational block. The
error radius is uniform in $a$ and $q$ and decays like $H/4^q$: the analytic tail is fully
discharged.
\end{thm}

\begin{rem}
Consequently a cofinal $1/32$-separation of the raw approximant from every integer suffices for
\#249. The corpus verifies the analogous kill \emph{unconditionally} at exactly two exponents,
$a=4$ and $a=6$
(\lean{actualLcmTailOrbit_four_and_six_notMem_int}{TotientActualLcmShortKill.lean:35}, \ev{Cert},
\scale{fixed}, via \lean{certifiedKill_diagonal_t16}{DiagonalPincerCertificates.lean:2832} and
\lean{certifiedKill_diagonal_t64}{DiagonalPincerCertificateT64Endpoint.lean:1928}). Nothing beyond
$a=6$ is proved. The depth $q$ is not free: the supply predicate pins $2q{+}1 =
\mathtt{oddGuardedCanonicalAdjacentSuffixDepth}(2^a)$, so per scale $a$ there is exactly one
admissible depth, not a search over depths.
\end{rem}

\begin{defn}[What would close it]
$\mathtt{PowerTwoActualLcmOrbitSeparationSupply}$: \ev{Open}, \scale{cofinal}.
\[
  \forall a_0\ \exists a\ge\max(2,a_0)\ \exists q,\
  \mathtt{oddGuardedCanonicalAdjacentSuffixDepth}(2^a)=2q{+}1 \wedge
\]
\[
  \forall z\in\mathbb Z,\ \tfrac1{32} + \mathtt{actualLcmRawErrorRadius}\ a\ q \;\le\;
  |\mathtt{actualLcmTailOrbit}\ a - z| .
\]
Since the error radius is explicit, this reduces to a distance-to-nearest-integer lower bound for
one explicitly computable rational per scale --- the cleanest reduction of the \#249 obligation to
a purely finite, computable question in the whole corpus.
\end{defn}

\subsubsection{SK-02 --- the short-window arithmetic kill, verbatim, truncated at $a_0\le 6$ (\coord{short-window-arithmetic})}

\begin{thm}
\lean{powerTwoActualLcmShortArithmeticKillSupply_through_six}{TotientActualLcmShortKill.lean:54}
and
\lean{irrational_totientSeries_of_shortArithmeticKillSupply}{TotientActualLcmShortKill.lean:63}.
\ev{Cert} for the witness, \ev{Lean} for the consumer. \scale{bounded}.
\[
  \forall a_0\le 6,\ \exists a,L,\ a_0\le a \wedge L < 2\cdot 2^a \wedge
  \mathtt{LcmDiagonalArithmeticKill}(2^a)\,L
\]
--- literally the open cofinal supply predicate below with its universal quantifier truncated at
$a_0\le 6$. This is the purest \scale{fixed}-vs-\scale{cofinal} row in the corpus: the proof term is
$\langle$\texttt{6, 93, ha0, by norm\_num, lcmDiagonalArithmeticKill\_two\_pow\_six}$\rangle$ --- a single
hard-coded witness $(a,L)=(6,93)$, itself discharged from
\lean{certifiedKill_diagonal_t64}{DiagonalPincerCertificateT64Endpoint.lean:1928} (a
\texttt{norm\_num} evaluation over an explicit table of $\varphi$ values). There is no argument in
$a$ whatsoever; removing the bound $a_0\le 6$ requires an entirely new proof, not a re-run.
\end{thm}

\begin{defn}[What would close it]
$\mathtt{PowerTwoActualLcmShortArithmeticKillSupply}$: \ev{Open}, \scale{cofinal}.
\[
  \forall a_0\ \exists a,L,\ a_0\le a \wedge L < 2\cdot 2^a \wedge
  \mathtt{LcmDiagonalArithmeticKill}(2^a)\,L .
\]
The short-window restriction $L<2\cdot 2^a$ buys extra structure --- every non-divisor offset in
that window is a bare prime power, by \texttt{eq\_prime\_pow\_of\_not\_dvd\_periodLcm} --- so this
is a better-equipped target than the raw SEP-02 supply even though it is formally a stronger
statement (it implies certified-kill directly, without the separation-and-round step).
\end{defn}

\subsubsection{b11 --- the diagonal pincer certificate bank (\coord{diagonal-pincer})}

\begin{thm}
\lean{certifiedKill_diagonal_all_imported}{DiagonalPincerCertificates.lean:2860}, with
\lean{diagonalPincerCertificateScales}{DiagonalPincerCertificates.lean:2842} and
\lean{diagonalPincerKillDepth}{DiagonalPincerCertificates.lean:2844}. \ev{Cert}. \scale{fixed}.
\[
  \forall t\in\{1,2,3,4,5,7,8,9,11,13,16,17\},\ \exists L,\
  \mathtt{certifiedKill}\ (\mathtt{periodLcm}\ t)\ (\mathtt{periodLcm}\ t)\ L
\]
at depths $\{6,5,7,7,9,14,15,14,21,22,23,26\}$ respectively, extended by the separate T19\ldots T64
modules to 28 historical values through $t=64$ (the $t=64$ endpoint is
\lean{certifiedKill_diagonal_t64}{DiagonalPincerCertificateT64Endpoint.lean:1928}).
The later aggregate theorem closes every scale $t\le82$ with no holes
(\lean{ErdosProblems.Skip.LadderT67.exists_diagonalKill_le_82}{ErdosProblems/Skip/LadderT67.lean:71264}).
Every witness
is a \texttt{norm\_num} evaluation over a hard-coded block of $\varphi$ values --- e.g.\ the $t=17$
certificate lists $\varphi(12252241),\ldots,\varphi(24504506)$ explicitly --- so nothing in the
proof is a function of $t$.
\end{thm}

\begin{rem}[A stronger signal in the depth table than the theorem states]
\lean{certifiedKill_depth_floor}{TotientTailPeriodKiller.lean:79} forces $2(2H_t{+}L{+}2)<2^L$,
i.e.\ $L \gtrsim \log_2(4\cdot\mathtt{periodLcm}\ t)$ at any certified depth. Comparing that floor
against $\mathtt{diagonalPincerKillDepth}$ shows every landed certificate fires within a small
additive constant of the theoretical minimum depth, at every tested scale --- a numerically
supported (\ev{Cert}, not \ev{Math}) anti-concentration observation, not a theorem.
\end{rem}

\begin{prop}[What would close it]
\[
  \exists C\ \forall t_0\ \exists t\ge t_0\ \exists L\le \log_2(4\cdot\mathtt{periodLcm}\ t)+C,\quad
  \mathtt{certifiedKill}\ (\mathtt{periodLcm}\ t)\ (\mathtt{periodLcm}\ t)\ L .
\]
This is an anti-concentration statement about the diagonal word's residue at the first admissible
depth --- exactly the shape a doubling-orbit equidistribution argument would produce.
\end{prop}

\subsubsection{a12 --- one window, sixteen periods, wrong-way quantifiers (\coord{h-uniform-certificate})}

\begin{thm}
\lean{certifiedKill_all_upto_sixteen}{CarrySurvivorExtinction.lean:574} and
\lean{totient_series_ne_rat_of_den_dvd_pow_two_mul_mersenne_upto_sixteen}{CertificateKernel.lean:18843}
(shallower sibling \lean{certifiedKill_all_small}{TotientTailPeriodKiller.lean:404}). \ev{Cert}.
\scale{fixed}.
\[
  \forall h\in[1,16],\quad \mathtt{certifiedKill}\ h\ 14\ 9
\]
--- a single position $N=14$ and a single depth $L=9$ that simultaneously certify every period $h$
up to $16$. Consequence: $S\ne r$ for every rational $r$ with $r.\mathrm{den}\mid 2^{14}(2^h{-}1)$,
$1\le h\le 16$. Proof is \texttt{decide} on a 9-bit window; no part of it is a function of $N$ or of
the $h$-range.
\end{thm}

\begin{rem}[Quantifier inversion --- the only $h$-uniform row in the corpus]
The 249-supply obligation is $\forall h\ \forall N_0\ \exists N\ge N_0\ \exists L$. This row proves
$\exists N\ \exists L\ \forall h\in[1,16]$: the quantifiers are inverted, and the $h$-range is
finite. The inversion helps in one direction (one $(N,L)$ covers a whole block of periods, more
than the obligation asks) and hurts in the other ($N$ is pinned at $14$, not cofinal). This is the
only row in the corpus with $h$-uniformity, and it is the reason a scale-uniform version would be
unusually strong.
\end{rem}

\begin{prop}[What would close it]
A scale-uniform version of the same shape: $\exists f:\mathbb N\to\mathbb N$ with $f(N)\to\infty$
such that $\forall N_0\ \exists N\ge N_0\ \exists L$ with $\mathtt{certifiedKill}\ h\ N\ L$ for all
$h\le f(N)$. This implies the obligation immediately (fix $h$, take $N_0$ large enough that
$f(N)\ge h$) and, unlike the obligation, is a statement about one window at a time rather than a
per-period search.
\end{prop}

\subsubsection{c3 --- the Farey denominator floor and its stalled growth law (\coord{farey-window})}

\begin{thm}
\lean{tsum_totient_div_pow_two_ne_ratCast_of_den_le_79639646646701375323355774875831053}{CertificateKernel.lean:18337},
built from the classical mediant lemma \lean{farey_gap}{GapFareyBound.lean:51} and the window
sharpness certificate \lean{gap_check_window_1_240_le_79639646646701375323355774875831053}{GapFareyBound.lean:176}
/ \lean{gap_check_window_1_240_first_failure}{GapFareyBound.lean:225}. \ev{Lean} for the general
mediant lemma; \ev{Cert} for the $K=240$ window instantiation. \scale{fixed}. If $S$ is rational,
its reduced denominator exceeds $7.9639646646701375323355774875831053\times 10^{34}$ (the exact
numeral in the declaration name). This is the strongest unconditional statement about $S$ in the
corpus, logically independent of the certificate-supply reduction.
\end{thm}

\begin{rem}[Why re-running the same window buys nothing]
The cofinal upgrade of this statement --- $S\ne p/q$ for every bound $q$, not just $q =
7.96\times10^{34}$ --- is literally $\mathrm{Irrational}(S)$, so the target coordinate is right.
What blocks promotion is proved on disk:
\lean{gap_check_window_1_240_first_failure}{GapFareyBound.lean:225} establishes that the $K=240$
window bound is \emph{sharp}, at exactly the mediant $b{+}d$, so re-running the same argument at
the same window buys nothing further. Each new $K$ needs (i) a freshly committed $2^K$-scale
totient residue $V_K$ and (ii) fresh continued-fraction convergents of $V_K'/2^K$, both hard-coded
numerals in the current proofs. The growth of the bound $b_K{+}d_K$ is governed by the convergent
denominators of the underlying constant.  Rationality would force eventual stalling, so unbounded
growth would prove \#249; however, no converse reduction or proved logical equivalence is known.
The sentence is therefore a diagnosis of why this fixed-window method reaches the original
difficulty, not an iff theorem.
\end{rem}

\begin{obs}[A closed sub-route]
\lean{ReducedDenominatorUnitGapCert}{PrimitiveWeightCertificate.lean:30} and
\lean{PrimitiveReducedDenominatorUnitGapSupply}{PrimitiveWeightCertificate.lean:47}, the corpus's
own attempt at strengthening this window refinement via a unit-modulo-odd-part criterion, prove a
hard ceiling in
\lean{reducedDenominatorUnitGapCert_prime_pow_iff}{PrimitiveWeightCertificate.lean:54} and
\lean{reducedDenominatorCandidateCount_prime_pow_le_one}{PrimitiveWeightCertificate.lean:111}
(\ev{Lean}, \scale{n/a}): at a prime-power denominator, the strengthening rescues at most \emph{one}
extra lattice point. That route is closed --- it cannot be the source of a growth law.
\end{obs}

\begin{prop}[What would close it]
A proved growth law for the window family: a function $g$ with $g(K)\to\infty$ and a proof that for
every $K$ the $(N{=}1,K)$ gap check passes for all $q\le g(K)$. Equivalently, a lower bound on the
convergent denominators of the totient-window constant, uniform in $K$.
\end{prop}

\subsubsection{TA-01/TA-02 --- the two-adic pulse block reaches the arc centre and fails by an exponential (\coord{two-adic-pulse})}

\begin{thm}
\lean{exists_prime_totient_twoAdic_pulse_divisors}{TotientTwoAdicPulseBlock.lean:54},
\lean{exists_prime_deltaTotient_twoAdic_pulseBlock}{TotientTwoAdicPulseBlock.lean:211},
\lean{windowDiscrepancy_modEq_half_of_twoAdic_pulse}{TotientTwoAdicPulseBlock.lean:262},
\lean{eventual_integral_tailDiff_has_cofinal_twoAdic_half_pulse}{TotientTwoAdicPulseBlock.lean:327}.
\ev{Lean}. \scale{uniform} (the theorem is unconditional and holds for \emph{every} $K$, $H$, $B$,
not merely cofinally many). For every $K\ge 2$, every $H>K$, and every bound $B$, there are primes
$p>B$ with a length-$(K{-}1)$ zero prefix and a terminal half-turn, giving
\[
  A_{H,p-K,K} \equiv 2^{K-1} \pmod{2^K}.
\]
Under eventual integrality this transfers to an integer $z$ with $(z:\mathbb R) =
\mathrm{totientTail}(p{+}H) - \mathrm{totientTail}(p)$ and $z\equiv 2^{K-1}\pmod{2^K}$.
\end{thm}

\begin{rem}[Lands at the exact centre of the arc, fails by an exponential --- and why]
This lands the residue at the exact centre of the certified-kill arc: $2^{K-1}$ is maximally far
from $0$ modulo $2^K$, precisely the coordinate the obligation wants. Taking $N{=}p{-}K$,
$h{=}H$, $L{=}K$ gives radius $N{+}h{+}L{+}2 = p{+}H{+}2$, so $\mathtt{certifiedKill}$ requires
$2^{K-1} > p{+}H{+}2$. But $z\equiv 2^{K-1}\pmod{2^K}$ forces $|z|\ge 2^{K-1}$, while the directed
tail bound forces $|z| < p{+}H{+}2$; and the construction's own congruence
$p\equiv 1{+}2^{K-1}\pmod{2^K}$ (equivalently $v_2(p{-}1)=K{-}1$) forces $p\ge 1{+}2^{K-1}$. So
$p{+}H{+}2 > 2^{K-1}$ always, and no choice of $p$ rescues it. The quantifier order is inverted:
the theorem gives, for fixed $K$, cofinally many \emph{large} $p$; the certificate needs a $p$
\emph{small} relative to $2^{K-1}$, impossible for a single letter since $|\varphi(N{+}H)-\varphi(N)|
< N{+}H$. Depth-promotion is not the issue: it has already been performed, generalizing an
exponent-2 construction to every $K$ uniformly by CRT-gluing over a
$\mathrm{Unit}\oplus\mathrm{Fin}(K{-}1)\oplus\mathrm{Fin}(K{-}1)$ family.
\end{rem}

\begin{prop}[What would close it]
The same half-turn residue realised by an \emph{accumulated} window rather than a single terminal
letter: cofinally many $(h,N,L)$ with $A_{h,N,L}\equiv 2^{L-1}\pmod{2^L}$ and
$2^{L-1} > N{+}h{+}L{+}2$. Only the weighted sum $\sum \Delta\varphi\cdot 2^{L-1-j}$ can carry
magnitude $2^L$; the per-letter pulse provably cannot. Concretely: a pulse-block construction whose
zero prefix and half-turn are imposed on the \emph{word}, not on one delta.
\end{prop}

\subsubsection{d-3a/d-3b --- cofinal jump positions, one unproved slack sign (\coord{lcm-jump-slack})}

\begin{thm}[Cofinal producer, fully proved]
\lean{exists_periodLcm_strict_jump_ge}{DiagonalFreshLossBridge.lean:2635}. \ev{Lean}.
\scale{cofinal} (proved, not open): $\forall t_0\ \exists t\ge t_0$ with
$\mathtt{periodLcm}\ t < \mathtt{periodLcm}(t{+}1)$, witnessed by $p{-}1$ for any prime $p>t_0$.
The power-of-two specialization
\lean{CanonicalAdjacentSuffixPowerTwoPostJumpSlackSupply}{DiagonalFreshLossBridge.lean:2693} shows
the positions $t=2^a{-}1$ already suffice, so no position search is needed at all.
\end{thm}

\begin{defn}[The slack scalar and the open supply]
\lean{canonicalAdjacentSuffixCentralSlack}{DiagonalFreshLossBridge.lean:2651}:
\[
  \mathtt{canonicalAdjacentSuffixCentralSlack}\ t = \min\bigl(d-2^{m-5},\ (2^m-2^{m-5})-d\bigr),
\]
with $m$ the canonical adjacent-suffix depth and $d$ the residue at that depth. The open
producers, both \ev{Open} \scale{cofinal}, are
\lean{CanonicalAdjacentSuffixPostJumpSlackSupply}{DiagonalFreshLossBridge.lean:2686} (any strict
jump) and \lean{CanonicalAdjacentSuffixPowerTwoPostJumpSlackSupply}{DiagonalFreshLossBridge.lean:2693}
(power-of-two positions), each asking $0\le \mathtt{canonicalAdjacentSuffixCentralSlack}(t{+}1)$
cofinally, and both compose to
$\mathrm{Irrational}(S)$ at
\lean{irrational_totientSeries_of_canonicalAdjacentSuffixPostJumpSlackSupply}{DiagonalFreshLossBridge.lean:2937}
and
\lean{irrational_totientSeries_of_canonicalAdjacentSuffixPowerTwoPostJumpSlackSupply}{DiagonalFreshLossBridge.lean:2944}.
\end{defn}

\begin{rem}[Census evidence only]
\ev{Cert}, \scale{fixed}. All $40$ strict jumps up to endpoint $113$ pass, closest margin
$\approx 0.000221$ of the modulus at $t=100$; the power-of-two endpoints $4,8,16,32$ pass with
slack fractions $0.41,0.036,0.40,0.19$. Census, not theorem. The structural handle nobody has used:
$\mathtt{periodLcm\_succ\_eq\_prime\_mul\_of\_strict\_jump}$
(\lean{periodLcm_succ_eq_prime_mul_of_strict_jump}{DiagonalFreshLossBridge.lean:2332}) says a
strict jump at $t$ means $t{+}1=p^k$ and $\mathtt{periodLcm}(t{+}1) = p\cdot\mathtt{periodLcm}(t)$
--- so the required statement is a transfer lemma for how the adjacent-suffix residue at height $H$
moves under $H\mapsto pH$.
\lean{canonicalAdjacentSuffixCentralSlack_eq_of_periodLcm_eq}{DiagonalFreshLossBridge.lean:2658}
already proves the slack is constant across each LCM plateau, so only jump indices matter.
\end{rem}

\begin{prop}[What would close it]
$0\le \mathtt{canonicalAdjacentSuffixCentralSlack}(2^a)$ for cofinally many $a$.
\end{prop}

\subsubsection{e2 --- the prime-jump commutator, one fixed witness (\coord{prime-jump-commutator})}

\begin{thm}
\lean{primeJumpSharpRadius}{PrimeJumpWindow.lean:37} and
\lean{primeJumpTailCommutator_notMem_int_of_sharpKill}{PrimeJumpWindow.lean:178}. \ev{Lean}.
\scale{uniform}. Uniform consumer for all $H,p,L$: a residue of the four-vertex commutator
$J(H,p)$ outside the sharp radius $3pH{+}(p{+}1)(L{+}2)$ forces $J(H,p)\notin\mathbb Z$. This is
tighter than the earlier $4pH$ two-cell disjunction.
\end{thm}

\begin{thm}[The one witness]
\lean{primeJumpSharpKill_twelve_five}{PrimeJumpWindow.lean:186}. \ev{Cert}. \scale{fixed}.
$\mathtt{primeJumpSharpKill}\ 12\ 5\ 15$ (i.e.\ $H=\mathtt{periodLcm}\ 4=12$, $p=5$, $L=15$),
proved by \texttt{decide} on an explicit 15-bit window with no dependence on $t$ or $p$; only a
single instance exists, at $t=4$, the smallest nontrivial height. The full endpoint composition is
\lean{irrational_totient_series_of_primeJumpSharpKill_supply}{PrimeJumpWindow.lean:193}.
\end{thm}

\begin{rem}
The supply hypothesis $\forall t_0\ \exists t\ge t_0\ \exists p,L{>}0$ with
$\mathtt{primeJumpSharpKill}(\mathtt{periodLcm}\ t)\ p\ L$ is a genuinely cheaper target than the
raw SEP-02 supply: it asks for one fresh prime $p$ per LCM height rather than a full central-arc
certificate. Because the consumer route is
\texttt{rational\_totient\_series\_forces\_lcm\_cone\_flatness}$\to$contradiction, the natural
attack is to pick $p$ as the fresh prime introduced at the next strict LCM jump (linking this row
to d-3a), where the commutator's old-channel contributions are annihilated by
\lean{oldChannel_affine_moment_annihilation}{JointExponentTransport.lean:36}.
\end{rem}

\subsubsection{b7 --- the M\"obius--Mersenne denominator channel: unbounded, but off-coordinate (\coord{mobius-mersenne})}

\begin{thm}
\lean{upperHalfMersenneProduct_lower_bound}{MersenneShadowDenominatorGrowth.lean:60},
\lean{lcmHeight_scaledMobiusShadow_den_lower_bound}{MersenneShadowDenominatorGrowth.lean:85},
\lean{exists_upperHalf_mersenne_channel_dvd_den_with_bounds}{MersenneShadowDenominatorGrowth.lean:99}.
\ev{Lean}. \scale{uniform}. For every $t\ge 5$, uniformly in $t$:
\[
  2^{t/2} \;\le\; \prod_{p\in \mathrm{upperHalfPrimes}(t)} \mathtt{mersenne}(p) \;\le\;
  \mathrm{den}\bigl(\mathtt{lcmHeight}(t)\cdot \mathtt{numericMobiusShadow}(\mathtt{lcmHeight}(t))\bigr).
\]
An unbounded denominator lower bound at every LCM height, with an individual surviving channel
isolated ($2^{t/2}\le \mathtt{mersenne}(p) < 2^t$). The proof is genuinely uniform in $t$
(Bertrand's postulate via \texttt{upperHalfPrimes\_nonempty}, plus $2^{p-1}\le 2^p{-}1$ factorwise
--- no table, no constant degrading with $t$).
\end{thm}

\begin{rem}[The corpus's only proved unbounded-growth quantity in this coordinate, but the wrong coordinate for \#249]
This lives in a different coordinate from the 249-supply obligation: it bounds the reduced
denominator of the \emph{scaled} shadow at LCM height $t$, and there is no landed transport from a
denominator lower bound to $\mathtt{certifiedKill}$ or to a totient-tail non-integrality. Two
further honest caveats are recorded in the module's own docstring: it does not rule out
cancellation by the foreign-defect term, and the growth rate $2^{t/2}$ is far below the scale
$\mathtt{lcmHeight}(t)\approx 2^{1.44t}$ against which the enclosure consumers measure error. The
existing consumers built on it
(\texttt{scaleFullTarget\_miss\_of\_lambert\_projected\_separation},
\texttt{scaleFullTarget\_miss\_of\_lambert\_projected\_num\_gap}) are \#257-facing
(\texttt{ScaleFullTargetHit}), not \#249-facing.
\end{rem}

\begin{prop}[What would close it]
Either (i) a Dirichlet-criterion bridge in the shape of
\lean{irrational_of_den_mul_abs_sub_tendsto_zero}{CertificateKernel.lean:5371}: a sequence of
rationals $u_t$ with $u_t\ne S$ and $\mathrm{den}(u_t)\cdot|S-u_t|\to 0$, which requires the
approximation error to beat the denominator, not merely the denominator to grow; or (ii) a
transport map from ``channel of size $\ge 2^{t/2}$ survives in the denominator'' to ``windowDiscrepancy
residue avoids the arc,'' which does not exist in the corpus. Route (i) needs the growth rate
raised from $2^{t/2}$ to beat the shadow's own truncation error.
\end{prop}

\subsubsection{d4/d5 --- the parallel rank obligation, and a corrected pair of citations (\coord{carry-rank})}

This row runs a second, independent open obligation in parallel to 249-supply, in the carry-kernel
\emph{rank} coordinate rather than the residue coordinate. \textbf{Correction}: the source index
attributed the two theorems below to swapped files; direct reads this session confirm the correct
attribution is as given.

\begin{thm}
\lean{linearIndependent_canonicalTotientKernelFamily}{TotientMahlerDefect.lean:900}. \ev{Lean}.
\scale{uniform}. Unconditional and uniform in $e$: the dyadic totient-kernel family is linearly
independent at every depth $e$ (dimension exactly $2^e{+}1$, via CRT and Dirichlet on primes in
arithmetic progression).
\end{thm}

\begin{thm}
\lean{not_irrational_totientSeries_implies_unbounded_carryRank_unconditional}{TotientCarryKernelRigidity.lean:284}.
\ev{Lean}. \scale{uniform}. If $S$ is rational there is a tempered integral carry orbit $u$ with
\[
  2^e - 1 \;\le\; \mathrm{finrank}_{\mathbb Q}\,\mathrm{span}(\mathtt{canonicalCarryKernelFamily}\ u\ e)
  \qquad\text{for every } e .
\]
\end{thm}

\begin{rem}[The scale side is finished; the missing direction is proved dead on one route]
So: rationality forces \emph{unbounded} carry-kernel rank, and the $2^e{-}1$ floor holds for all
$e$ with a uniform proof. What is missing is the opposite inequality --- a rationality-side rank
upper bound, which would contradict the floor and close \#249. The corpus proves one natural route
to it is dead:
\lean{false_of_compressedAdjointCertificate}{TotientMahlerDefect.lean:1051} (\ev{Lean},
\scale{n/a}) shows no
\lean{CompressedAdjointCertificate}{TotientMahlerDefect.lean:1039} ($Q\cdot v\cdot A = $ boundary
with $|\mathrm{boundary}| < Q\cdot v$ and $A\ne 0$) can exist, and
\texttt{not\_finiteDimensional\_span\_fullTotientKernel} (immediately above it in the same file)
shows the full family spans an infinite-dimensional space. Cross-checked against a separate row,
$\mathtt{lcm\_factorIdeal\_finiteRank\_shiftAlgebra\_not\_sufficient}$: a single explicit
countermodel defeats every finite-rank shift-polynomial observation simultaneously, so finite-rank
amplification per se is not the missing rigidity.
\end{rem}

\begin{prop}[What would close it]
A rationality-side rank upper bound: $\exists C$ such that any tempered integral binary carry
orbit for a rational $\mathtt{binaryCoeffSeries}$ has
$\mathrm{finrank}_{\mathbb Q}\,\mathrm{span}(\mathtt{canonicalCarryKernelFamily}\ u\ e)\le C$ for
all $e$, or any bound growing slower than $2^e{-}1$.
\texttt{not\_irrational\_totientSeries\_implies\_mod\_period\_and\_unbounded\_rank}
(\texttt{TotientTailCarryPeriod.lean:224}) pins the obstacle precisely: rationality buys uniform
eventual periodicity of $u$'s dyadic sections mod $v$, and that periodicity provably does
\emph{not} promote to a $\mathbb Q$-rank bound without extra arithmetic input.
\end{prop}

\subsection{Synthesis: where the frontier actually is}

Collecting the fourteen rows, three facts stand out as loci for future effort, stated with the same
precision as the rows themselves rather than as summary rhetoric.

First, the single closest approach to 249-supply is the first-harmonic reduction of
\S\ref{ssec:headline}: a constant-saving ($21/25$ or $9/10$) exponential-sum cancellation estimate,
with a satisfiable room condition and, via the subset consumer
(Theorem~\ref{thm:hgap-subset}), no requirement that the saving hold on a full interval or a
positive-density subset. This is the only row in the index whose remaining content is recognisably
a single classical estimate rather than a compound arithmetic-geometric statement.

Second, the top-edge staircase (TE-04 through TE-05-weakest, \coord{actual-lcm-top-edge}) is the
only row where the corpus has both halved the obligation (one-sided instead of symmetric) and
exposed, via the exact identity of Observation~\ref{obs:staircase-tension}, precisely which of two
logically symmetric escape branches survives the sign machinery already proved elsewhere in the
same file family. That five strictly weaker equivalent or sufficient forms are already proved
inter-derivable (Proposition~\ref{prop:te-chain}) means effort here is not fragmented across
restatements: proving any one closes the others automatically.

Third, the d4/d5 rank obligation is structurally independent of the residue-coordinate rows above:
it is a second, self-contained sufficient condition for \#249 (an upper rank bound contradicting
the proved $2^e{-}1$ lower bound), in a coordinate (carry-kernel rank) where a natural finite-rank
strengthening is already proved impossible. A rank upper bound, if found, would not need to route
through $\mathtt{certifiedKill}$ at all.

No row in this index is a solution, a partial solution in the sense of resolved cases, or evidence
that \#249 is likely true or false. Every open producer listed is exactly as open as the original
Erd\H{o}s--Borwein question it reduces to; what the index adds is the exact remaining content,
verified against the live Lean tree, of each reduction.
